TBM excavation is a moving spatial relation problem. As the cutterhead and shield advance, geological units enter, approach and leave the potential relation zone of different TBM components. Changes in shield pressure, displacement, thrust, torque, penetration and advance rate therefore reflect not only operational control but also the evolving spatial configuration between geological anomalies and machine geometry. The GIScience problem addressed here is the representation of these changing rock--machine relations under heterogeneous spatial supports.

\begin{figure}[t]
\centering
\includegraphics[width=0.96\textwidth]{fig_research_scene.pdf}
\caption{Research scene and object of representation. A TSP-derived geological voxel field and a component-labelled advancing TBM surface are observed under different spatial supports. The proposed representation constructs candidate rock--TBM relations between anomalous rock voxels and machine components, while monitoring responses are aligned to the resulting component-level indices for event interpretation.}
\label{fig:research_scene}
\end{figure}

In this paper, the engineering event is referred to as \emph{TBM shield jamming}. The term denotes a jammed state in which surrounding ground convergence and shield-side resistance make continued advance difficult or impossible under the available thrust. The term is preferred to ``shield sticking'' because the tunnelling literature more commonly uses TBM jamming or shield jamming when describing squeezing-ground obstruction, shield skin-friction resistance, and jamming accidents. ``Stuck'' or ``jammed'' is used only as a state descriptor, not as the main technical term.

Monitoring-based TBM response models are useful because short-horizon excavation responses often have strong temporal persistence. Empirical, machine-learning and sequence-based models can learn response inertia from rock, machine and recent operational histories \citep{yagiz2008utilizing,hassanpour2010tbm,armaghani2018prediction,hochreiter1997long}. However, a monitoring curve does not state which geological voxel is adjacent to the cutterhead or shield, whether that adjacency is geometrically plausible, or how the relation support changes as excavation proceeds. This limits the spatial interpretability of monitoring-only analyses.

Advance geological prediction and TSP surveys provide a complementary view by describing spatial heterogeneity ahead of or around the excavation face \citep{li2019tunnel}. These data can be represented as chainage-referenced geological voxels or attribute fields, preserving where low-velocity or other anomalous ground occurs. Yet a geological voxel field answers where an anomaly is located, not which TBM component is a plausible relation partner at a particular excavation step. The missing object is therefore neither another monitoring feature nor a single geological average, but a component-indexed relation support between geological spatial units and machine-side spatial units.

A moving buffer or a linear-referenced overlay is insufficient for this task. A buffer can summarise the mean anomaly inside a moving active zone, but it does not retain which rock voxel is related to which TBM component, which edge attributes make the relation plausible, or whether the relation appears, persists or disappears over excavation steps. The proposed graph keeps heterogeneous entities, component labels, candidate-edge attributes and edge lifecycle explicit. The resulting component--excavation-step index is a relation aggregation over screened candidate edges rather than a simple voxel-field average.

This paper therefore treats rock--TBM interpretation as a spatial entity--relation modelling problem. The input consists of a TSP-derived rock voxel field, an advancing parameterized TBM surface and excavation-step-aligned monitoring records. Two graph entity types are defined: rock voxels and TBM surface nodes. Monitoring responses are aligned to graph outputs but are not used to define graph nodes or rock--TBM edges. The method constructs a dynamic rock--TBM spatial relation graph at each excavation step using active-zone, distance, surface-normal compatibility and excavation-state constraints. Weighted candidate edges are aggregated by TBM component to compute a component-level geological anomaly index.

The paper makes three contributions. First, it defines a rock--TBM spatial entity model that places TSP-derived rock voxels and component-labelled TBM surface nodes in a common excavation-step coordinate while keeping monitoring responses as aligned records. Second, it constructs a sparse per-step candidate relation graph using physically interpretable spatial constraints rather than full connectivity or unconstrained learned edges. Third, it computes and evaluates a component-level geological anomaly index that supports relation reconstruction, component specificity, comparison with alternative summaries, sensitivity checks and edge-level provenance. The empirical evidence is used to validate the spatial relation representation in a confirmed TBM shield-jamming event.
